% Nová příprava: 2026-10-28; 1. ročník SPŠ, etapa 13–24.
% Rozpis: 1_rocnik/tematicke_plany/Rozpis_vyuky_MAT_1_rocnik.tex, hodina 16.
% Původ: vlastní formulace výkladu a vlastní úlohy; metodická inspirace uvedenými zdroji.
% Kontrola, obsahové porovnání zdrojů a časy: planovani/DOKONCENI_PRIPRAV_13_24.md.
% Učebnice: ucebnice/01_Zakladni_poznatky/Mnoziny/Množiny.tex
% oddíly Rovnost množin, Podmnožiny, Množinové operace a Vennovy diagramy; dle tématu hodiny.
% Realisticky.cz: externi_zdroje/realisticky/01_Zakladni_poznatky/03_Mnoziny/010302_Mnoziny_II_P/010302_Lekce.pdf
% Realisticky.cz: externi_zdroje/realisticky/01_Zakladni_poznatky/03_Mnoziny/010303_Mnozinove_operace/010303_Lekce.pdf
% Realisticky.cz: externi_zdroje/realisticky/01_Zakladni_poznatky/03_Mnoziny/010304_Vennovy_diagramy/010304_Lekce.pdf
\documentclass[11pt,a4paper]{article}
% Spolecny styl pro ucitelske pripravy MAT (Work/MAT)
% Verze 1.1, 2026-10-23
\NeedsTeXFormat{LaTeX2e}
\ProvidesPackage{priprava}[2026/10/23 v1.1 Sablona ucitelskych priprav MAT]

% Zakladni typografie a cestina
\RequirePackage[T1]{fontenc}
\RequirePackage[utf8]{inputenc}
\RequirePackage[czech]{babel}
\RequirePackage{lmodern}
\RequirePackage{microtype}
\RequirePackage{amsmath,amssymb}
\RequirePackage{xcolor}
\RequirePackage{enumitem}
\RequirePackage[most]{tcolorbox}
\RequirePackage{tabularx,booktabs,array}
\RequirePackage{fancyhdr}
\RequirePackage{needspace}
\RequirePackage{geometry}
\RequirePackage[hidelinks]{hyperref}

\geometry{a4paper,top=17mm,bottom=18mm,left=23mm,right=23mm}
\definecolor{notegray}{gray}{0.39}
\setlength{\parindent}{0pt}
\setlength{\parskip}{2.5pt plus 1pt minus 0.5pt}
\setlist[itemize]{topsep=2pt,itemsep=1pt,parsep=0pt,leftmargin=1.45em}
\setlist[enumerate]{topsep=2pt,itemsep=2pt,parsep=0pt,leftmargin=1.65em}
\setcounter{secnumdepth}{0}
\renewcommand{\arraystretch}{1.12}

% Zahlavi: vlevo rocnik | globalni tema, vpravo cislo: nazev hodiny
% Zapati zustava zcela prazdne.
\pagestyle{fancy}
\fancyhf{}
\renewcommand{\headrulewidth}{0.35pt}
\renewcommand{\footrulewidth}{0pt}
\setlength{\headheight}{14pt}
\newcommand{\nastavPripravu}[4]{%
  \fancyhead[L]{\small #1\enspace\textbar\enspace #2}%
  \fancyhead[R]{\small #3: #4}%
}

% Minimalisticke cernobile boxy. Kazda definice zvlast.
\tcbset{commonlesson/.style={enhanced,breakable=false,colback=white,colframe=black,
  colbacktitle=white,coltitle=black,boxrule=0.45pt,
  arc=0.8mm,outer arc=0.8mm,left=3mm,right=3mm,top=1.6mm,bottom=1.6mm,
  toptitle=1mm,bottomtitle=0.7mm,fonttitle=\bfseries,
  attach boxed title to top left={xshift=2.5mm,yshift=-1.15mm},
  boxed title style={colback=white,colframe=white,boxrule=0pt,left=1mm,right=1mm}}}
\newtcolorbox{zadaniBox}[1][]{commonlesson,title={Zadání},#1}
\newtcolorbox{prikladBox}[1][]{commonlesson,title={Příklad},#1}
\newtcolorbox{definiceBox}[1][]{commonlesson,title={Definice},#1}

% Spolecny box pro zadani a jeho reseni, bez mezery mezi castmi.
% Pouziti:
%   \begin{prikladReseniBox}
%     Zadani ulohy.
%     \tcblower
%     \begin{align*} ... \end{align*}
%   \end{prikladReseniBox}
% Spodni cast ma jemny sedy podklad a delici linku; lze v ni sazet
% vice radku matematiky. Obe casti maji jen jeden spolecny ramecek.
\tcbset{solutionlesson/.style={
  commonlesson,
  enhanced jigsaw,
  breakable,
  colbacklower=black!3,
  segmentation style={solid,draw=black!25,line width=0.35pt},
  fontlower=\small\color{notegray},
  before lower={\noindent\textbf{Řešení:}\enspace},
  leftlower=3mm,rightlower=3mm
}}
\newtcolorbox{prikladReseniBox}[1][]{solutionlesson,title={Příklad},#1}
\newtcolorbox{zadaniReseniBox}[1][]{solutionlesson,title={Zadání},#1}

% Pomucky pro autora pripravy: sedive poznamky; \reseni pro starsi pripravy
\newcommand{\poznamka}[1]{\par\begingroup\small\color{notegray}\itshape
  \textnormal{\textbf{Poznámka pro učitele.}}\enspace #1\par\endgroup}
\newcommand{\reseni}[1]{\par\begingroup\small\color{notegray}
  \textbf{Očekávané řešení:}\enspace #1\par\endgroup}
\newcommand{\blok}[1]{\needspace{5\baselineskip}\subsection*{#1}}
\newcommand{\vyrok}[1]{\uv{#1}}

\usepackage{tikz}
\usetikzlibrary{patterns,calc,angles,quotes}
\nastavPripravu{1. ročník}{Množiny}{16}{Rozdíl a doplněk}
\begin{document}
\blok{Vybrat prvky, které nesplňují druhou podmínku}
Z dílů určených ke kontrole chceme ponechat ty, které ještě nebyly zkontrolovány.
Musíme říci, odkud vybíráme, i co vylučujeme. Připomeň průnik a sjednocení,
ale jejich definice dnes znovu nevykládej.

\begin{definiceBox}
\textbf{Rozdíl} $A\setminus B$ tvoří právě prvky $A$, které nepatří do $B$:
$x\in A\setminus B\Leftrightarrow (x\in A\land x\notin B)$.
\end{definiceBox}
$B$ nemusí být podmnožinou $A$. Prvky $B$, které v $A$ vůbec nejsou,
výsledek neovlivňují. Na pořadí obou množin záleží.
\begin{prikladReseniBox}[title={P1: dva různé rozdíly},breakable=false]
Nechť $A=\{-2,0,2,4\}$ a $B=\{0,1,4,6\}$. Urči oba rozdíly.
\tcblower
$A\setminus B=\{-2,2\}$ a $B\setminus A=\{1,6\}$.
Společné prvky $0,4$ nejsou ani v jednom rozdílu. Výsledky se nerovnají.
\end{prikladReseniBox}

\begin{definiceBox}
Pro $A\subseteq U$ je \textbf{doplněk množiny $A$ vzhledem k základní množině $U$}
množina $A'_U=U\setminus A$. Je-li $U$ jasně určeno, píšeme také $A'$.
\end{definiceBox}
Bez základní množiny není doplněk určen jednoznačně. Doplněk není
„soubor všeho na světě mimo $A$“. Rozdíl $A\setminus B$ je doplňkem
$A\cap B$ v $A$; jako doplněk samotného $B$ v $A$ jej zde pojmenujeme jen při $B\subseteq A$.

\begin{center}
\begin{tikzpicture}[scale=.72]
\begin{scope}
\clip (-.8,0) circle(1.3);
\fill[pattern=north east lines] (-2.2,-1.5) rectangle (2.2,1.5);
\fill[white] (.8,0) circle(1.3);
\end{scope}
\draw (-.8,0) circle(1.3); \draw (.8,0) circle(1.3);
\node at (-1.3,1.5) {$A$}; \node at (1.3,1.5) {$B$};
\node at (0,-1.8) {$A\setminus B$};
\begin{scope}[xshift=6.7cm]
\fill[black!15] (-2.3,-1.5) rectangle (2.3,1.6);
\fill[white] (0,0) circle(1.1);
\draw (-2.3,-1.5) rectangle (2.3,1.6);
\draw (0,0) circle(1.1);
\node at (0,0) {$A$}; \node at (2,1.3) {$U$};
\node at (0,-1.8) {$A'_U=U\setminus A$};
\end{scope}
\end{tikzpicture}
\end{center}

\begin{prikladReseniBox}[title={P2: stejné $A$, jiné $U$},breakable=false]
Pro $A=\{2,4\}$ urči doplněk nejprve v $U_1=\{1,2,3,4,5\}$,
potom v $U_2=\{0,1,2,3,4,5,6\}$.
\tcblower
$A'_{U_1}=\{1,3,5\}$, ale $A'_{U_2}=\{0,1,3,5,6\}$.
Obě základní množiny $A$ obsahují, výsledek závisí na zvoleném okolí.
\end{prikladReseniBox}
\begin{prikladReseniBox}[title={P3: nekonečná základní množina},breakable=false]
Urči doplněk $\mathbb N$ v $\mathbb Z$ a doplněk intervalu
$I=\langle-1;3)$ v $\mathbb R$.
\tcblower
$\mathbb N'_{\mathbb Z}=\{x\in\mathbb Z\mid x\le0\}$, protože naše $\mathbb N$ začíná jedničkou.

$I'_{\mathbb R}=(-\infty;-1)\cup\langle3;+\infty)$.
Bod $-1$ v $I$ je, proto v doplňku není; bod $3$ naopak do doplňku patří.
\end{prikladReseniBox}
Přímo z definice: $A\setminus A=\varnothing$, $A\setminus\varnothing=A$,
$A\cap A'_U=\varnothing$, $A\cup A'_U=U$ a $(A'_U)'_U=A$.
V posledním vztahu musí zůstat stejné $U$.

\blok{Samostatné procvičení}
\begin{zadaniReseniBox}[title={S1: rozdíl, doplněk a krajní body},breakable=false]
Použij $U=\{0,1,2,3,4,5,6,7\}$, $A=\{0,2,4,6\}$, $B=\{2,3,6\}$.
\begin{enumerate}[label=\alph*)]
\item Urči $A\setminus B$, $B\setminus A$, $A'_U$, $B'_U$.
\item Urči $(A\cup B)'_U$ a $(A\cap B)'_U$.
\item Urči $\varnothing'_U$ a $U'_U$ a vysvětli rozdíl mezi těmito výsledky.
\item V $\mathbb R$ urči doplněk $(-\infty;2\rangle$ a doplněk $\{0\}$.
\end{enumerate}
\tcblower
a) $\{0,4\}$, $\{3\}$, $\{1,3,5,7\}$, $\{0,1,4,5,7\}$.
b) $A\cup B=\{0,2,3,4,6\}$, doplněk $\{1,5,7\}$;
$A\cap B=\{2,6\}$, doplněk $\{0,1,3,4,5,7\}$.
c) $\varnothing'_U=U$ a $U'_U=\varnothing$: v prvním případě nic nevyloučíme,
ve druhém vyloučíme všechny prvky základní množiny.
d) $(2;+\infty)$ a $(-\infty;0)\cup(0;+\infty)$; krajní body jsou vyloučené.
\end{zadaniReseniBox}
\poznamka{Jádro: P1–P3, význam pěti jednoduchých vztahů a S1,
přibližně 38–39 min práce, zbytek organizace a závěr.
U S1 b) zatím jen spočti konkrétní výčty; obecné zákony přijdou v hodině 17.}

\blok{Volitelná rezerva}
\begin{zadaniReseniBox}[title={R1: právě jedna podmínka},breakable=false]
Pro množiny ze S1 urči $(A\setminus B)\cup(B\setminus A)$.
Slovně ji popiš. Může být $A\setminus B=B\setminus A$ neprázdná množina?
\tcblower
Výsledek je $\{0,3,4\}$: prvky patřící právě do jedné ze dvou množin.
Oba rozdíly nemohou mít společný prvek: první vyžaduje příslušnost do $A$,
druhý ji vylučuje. Jsou-li si rovny, musí být oba prázdné; pak $A=B$.
\end{zadaniReseniBox}
\poznamka{Rezerva 8–10 min. Není nutné zavádět nový symbol pro symetrický rozdíl.}

\blok{Domácí procvičení}
\begin{zadaniReseniBox}[title={D1: 12–18 minut}]
V a) a b) je $U=\{-2,-1,0,1,2,3,4\}$, $K=\{-2,0,2,4\}$, $L=\{-1,0,1,2\}$.
\begin{enumerate}[label=\alph*)]
\item Urči $K\setminus L$, $L\setminus K$, $K'_U$ a $L'_U$.
\item Urči $(K\setminus L)'_U$.
\item V $\mathbb Z$ urči doplněk $\{x\in\mathbb Z\mid x<3\}$.
\item V $\mathbb R$ urči doplněk $(0;5\rangle$ a ověř příslušnost $0$ a $5$.
\item Proč nelze bez další informace určit doplněk $\{1,2\}$?
\end{enumerate}
\tcblower
a) $\{-2,4\}$, $\{-1,1\}$, $\{-1,1,3\}$, $\{-2,3,4\}$.
b) $\{-1,0,1,2,3\}$.
c) $\{x\in\mathbb Z\mid x\ge3\}$.
d) $(-\infty;0\rangle\cup(5;+\infty)$; nula patří do doplňku, pětka ne.
e) Chybí základní množina. V $U=\{1,2\}$ je doplněk prázdný,
v $U=\{1,2,3\}$ je $\{3\}$.
\end{zadaniReseniBox}
\end{document}
